\paragraph{Correlated instantaneous frequency.}
A separate prescribed drive model replaces Wiener phase increments by
$\theta(t)=\theta_0+\int_\Delta^t\xi(s)\,ds$, with an independent uniform
initial angle and stationary zero-mean Ornstein--Uhlenbeck frequency
$\mathbb E\xi(t)\xi(s)=(D/\tau_c)e^{-|t-s|/\tau_c}$.
Its phase correlation at lag $u\ge0$ is
\[
 C(u)=\exp\{-D[u-\tau_c(1-e^{-u/\tau_c})]\}.
\]
This retains the pathwise source norm and fixes the long-time phase
diffusion constant $D$; the instantaneous frequency variance is
$D/\tau_c$. The Wiener kernel is recovered as $\tau_c\to0$ at fixed $D$.
Finite $\tau_c$ removes its linear short-lag cusp, but does not impose a
hard bandwidth cutoff or define a thermal bath. Pointwise larger phase
correlation need not give larger work against an oscillatory response.

\paragraph{Small correlation-time limit.}
For a fixed finite spatial operator and fixed $D\geq0$, let $\tau=\tau_c$ and write the exact mean work as
$m_\tau=\int_0^T W(u)C_\tau(u)\,du$, with real lag kernel $W$ and
$m_D=\int_0^T W(u)e^{-Du}\,du$. If $|W(u)|\leq M$, then
\[
 |m_\tau-e^{D\tau}m_D|\leq
 \frac{e^{D\tau}DM\tau^2}{1+D\tau},\qquad m'_+(0)=Dm_D.
\]
If $W$ is continuous at zero, the present source-work observable has
$W(0)=2B_{\mathrm{src}}$, where $B_{\mathrm{src}}=\int_0^T\|S(t)\|^2dt$ is the exact
nominal source budget. Consequently,
$m_\tau=m_D+Dm_D\tau+(D^2m_D/2-2DB_{\mathrm{src}})\tau^2+o(\tau^2)$.
These local asymptotics neither determine the comparison at $\tau=0.1$
nor imply global monotonicity, a thermal interpretation, or a continuum error bound.

For the fixed choice $D=1$, $\tau_c=1/10$, put $a=1/10$.
The mean source work has the signed representation
$m_{\mathrm{OU}}=e^a\sum_{n\ge0}(-a)^n m_{1+10n}/n!$, where $m_\lambda$
is the two-quadrature work obtained from the auxiliary generator
$A-\lambda I$. This is not a positive mixture of dissipative states.
The selected truncation retains $n=0,\ldots,8$. If
$B\ge2\int_{t>s}|K_{\mathrm{work}}(t,s)|\,dt\,ds$, its analytical work-tail
bound is
\[
 \epsilon_{\mathrm{tail}}\le
 B\,\frac{10}{9}\,\frac{(1/10)^9}{9!}\,\frac{1}{1-1/100}.
\]
The finite-model bound uses a positive reference energy norm,
a coefficient-envelope check and an exact-rational source-norm bound
for frozen nominal binary coefficients. This bounds the omitted series,
not floating exponential evaluation, signed-sum roundoff or RK integration.
Those errors and the combined space/time sensitivity require separate
assessment; auxiliary endpoint energies are not ensemble energies.
The first finite-correlation implementation passed its scalar controls but
failed the preselected modified quadratic-balance tolerance on the coarser
field grid. That implementation therefore yielded no accepted finite-correlation
field result, and its finer grid was not run. The analytical series-tail estimate
did not repair this numerical acceptance failure, and the original outcome
and thresholds remain unchanged. A separate calculation with a different
integrator and a prospectively specified temporal-sensitivity comparison is
reported below.

\paragraph{A separate finite-correlation pulse comparison.}
A separately specified two-stage Gauss calculation retained the prescribed
source, its normalization, $D=1$, $\tau_c=0.1$ and the signed nine-term expansion.
It split integration steps at the two known source-window joins and used
both genuine quadratures for every auxiliary rate, plus a zero-source
control. Three fixed integrations compared $(h,\Delta t)=(0.04,0.005)$,
$(0.04,0.0025)$ and $(0.02,0.0025)$ through $T=3$ at $R=20$.
All 108 recorded acceptance checks passed; this comparison was made against
the previously stored Wiener-source work on the corresponding spatial grid,
not against an auxiliary endpoint energy or a newly repeated Wiener run.

On the finer spatial grid, the truncated OU mean source work was
$2.1932622828\times10^{-4}$, compared with $2.1733143747\times10^{-4}$
for the stored Wiener result. The difference, $1.9947908110\times10^{-6}$
(approximately $0.918\%$), exceeded the predeclared diagnostic guard
$1.2658579115\times10^{-8}$ and had the same positive sign on the coarser grid.
The guard includes the analytical series-tail allowance and measured temporal,
spatial, adapter and accounting contributions; it is not a certified total-error
bound. The fixed coarse-grid time comparison does not directly measure fine-grid
temporal error. This is a bounded increase in mean source work for one prescribed
drive in the frozen finite-domain model, not evidence of retained energy,
heating, formation, improved efficiency or experimental realizability.

\begin{figure}[htbp]
\centering
\includegraphics[width=.98\linewidth]{mean-injected-work.pdf}
\caption{Cumulative mean source work for prescribed OU and Wiener phase
drives on the $h=0.02$ grid (left), and their difference (right).
The OU curve is the signed nine-term mean of the two genuine quadrature
works per auxiliary rate; the Wiener curve uses the previously stored
calculation. Both use integrated original source work, not auxiliary
endpoint or ensemble energy. Straight lines join 31 stored times, with no
additional field integration for the figure. The endpoint marker spans
the predeclared diagnostic guard; it is neither a confidence interval nor
a certified total-error bar. All quantities are in dimensionless model units.}
\label{fig:ou-source-work}
\end{figure}

\paragraph{One exterior-domain sensitivity check.}
A separate calculation at fixed $h=0.04$, $\Delta t=0.0025$ and $T=3$
extended $R=20$ to $R=30$, retaining byte-identical source values and
interior coefficients and comparing with the stored $R=20$ calculation.
Across the 31 shared sample times, the absolute-weighted majorant for the
change in the same-method OU-minus-$n=0$ source-work contrast was
$1.14\times10^{-19}$, below the predefined diagnostic threshold
$1.27\times10^{-9}$. This is agreement at floating-point roundoff scale,
not absolute accuracy at $10^{-19}$; common spatial, temporal and profile
errors can remain. The check does not establish an infinite-domain limit,
long-time behavior or control between samples, and neither replaces the
earlier Wiener baseline nor reduces the original comparison guard.

\paragraph{Signed contributions from source-time separations.}
A separate response-kernel calculation uses the fixed $h=0.04$, $R=20$
spatial operator and both genuine quadratures of each source profile.
After averaging the uniform initial phase, its mean source work is
$m[C]=\int_0^1 W(u)C(u)\,du$, where $u$ is the separation of two source
times, not the observation time. The kernel can have either sign, with
$W(0)=2B_{\mathrm{src}}$. On the finest of three lag meshes
($400$, $800$, $1600$ subdivisions), with independent overlap quadratures
of orders $16$ and $32$, the four prescribed quarter-interval contributions
to $\int_0^1 W(u)[C_{\mathrm{OU}}(u)-e^{-u}]\,du$ are, in units of $10^{-6}$,
\[
 (+23.13074873,\;-9.08245885,\;-11.01641674,\;-1.03681136).
\]
Their signed sum is $1.99506177879\times10^{-6}$, with summed heuristic
refinement envelopes $6.98\times10^{-14}$; these are not a certified
error bound or confidence interval. The shortest-lag contribution exceeds
the three negative later-lag contributions (Fig.~\ref{fig:signed-lag-kernel}).
Closure to the stored same-grid Gauss OU and Wiener-rate works passes the
predeclared diagnostic allowances. The separate shift of the Gauss
Wiener-rate work relative to the old RK Wiener baseline is
$-8.47524\times10^{-10}$; it is not assigned to any lag bin.
This calculation is not the $h=0.02$ comparison above and does not exactly
decompose its mixed-integrator result. A lag bin integrates pairs of source
times, not instantaneous positive or negative power. All separations are
at most one model time unit, below the nominal period $2\pi/\rho\simeq3.44$;
the response involves the full generator, not an isolated BIC or a long
ringdown. No stored-energy, formation or stabilization conclusion follows.

\begin{figure}[htbp]
\centering
\includegraphics[width=.98\linewidth]{signed-lag-kernel.pdf}
\caption{Stored finest-grid signed lag kernel $W(u)$ (left) and four fixed
integrals of $W(u)[C_{\mathrm{OU}}(u)-e^{-u}]$ (right), for the frozen $h=0.04$
model. Lines connect the stored lag nodes without smoothing. The caps and
reported net envelope are heuristic numerical sensitivity measures, too
small to resolve at this scale, not confidence intervals or certified
errors; no continuous kernel error band is asserted. The positive net
source-work difference retains all three negative contributions. These
are source-time-pair contributions, not a time history of instantaneous
power or retained energy.}
\label{fig:signed-lag-kernel}
\end{figure}
